\chapter{Soluciones de los simulacros adicionales C a E}
Las cifras se calculan con las operaciones mostradas; son aritmética reproducida de datos docentes, sin aprobación académica independiente.

\section{Simulacro C}
\textbf{Cuestiones.} El trabajo de Euler es el cambio del flujo de momento angular y solo depende de la cinemática, $w=U\,\Delta C_\theta$; la eficiencia determina qué parte de ese trabajo se convierte en aumento de presión total y qué parte se pierde como entropía, de modo que a igual $w$ una eficiencia menor da menor razón de presiones. El cociente $W_2/W_1$ mide la desaceleración relativa en el rotor y es un indicador de difusión: por debajo de unos $0.7$ el riesgo de desprendimiento aumenta, aunque el umbral depende del perfil y no se universaliza. El deslizamiento reduce el remolino de salida porque el flujo no sigue el álabe, con lo que cae $C_{\theta2}$ y el trabajo.

\textbf{Escalón axial.} $w=U(C_{\theta2}-C_{\theta1})=27500\,\mathrm{J/kg}$ y $\Delta T_0=27.3632\,\mathrm K$, de modo que $T_{02}=315.513\,\mathrm K$. La razón de presiones es $1.32432$ y $p_{02}=134187\,\mathrm{Pa}$. La potencia es $\dot m w=550000\,\mathrm W$ y la reacción $1-(C_{\theta1}+C_{\theta2})/(2U)=0.5$. Las velocidades relativas son $W_1=240.832\,\mathrm{m/s}$ y $W_2=174.642\,\mathrm{m/s}$, con ángulos $48.3665^\circ$ y $23.6294^\circ$ respecto al eje. El cociente $W_2/W_1=0.725163$ está justo por encima de $0.72$, valor que suele citarse como límite para evitar difusión excesiva; se trata de un criterio orientativo.

\textbf{Rodete.} $U_2=565.487\,\mathrm{m/s}$ y $C_{\theta2}=\sigma U_2=480.664\,\mathrm{m/s}$. El trabajo es $w=U_2C_{\theta2}=271809\,\mathrm{J/kg}$ y $\Delta T_0=270.457\,\mathrm K$, con razón de presiones $6.68999$. El Mach periférico es $U_2/a_{01}=1.62861$ con $a_{01}=347.221\,\mathrm{m/s}$, y la potencia con $1.2\,\mathrm{kg/s}$ es $326171\,\mathrm W$. Un Mach periférico supersónico obliga a vigilar el choque en el difusor, y la tensión del rodete limita la velocidad periférica; ninguno de los dos efectos entra en este cálculo ideal.

\subsection*{Procedimiento y operaciones del simulacro C}
Se aplica la ecuación de Euler con radio constante en el escalón axial y con $C_{\theta1}=0$ en el rodete. La razón de presiones sale del rendimiento total-total; la reacción, de la cinemática con $C_x$ constante.

\begin{enumerate}

\item Temperatura total de entrada.

\[\fitmath{T_{01}=288.15\approx 288.15\,\mathrm{K}}\]

\item Presión total de entrada.

\[\fitmath{p_{01}=101325\approx 101325\,\mathrm{Pa}}\]

\item Velocidad periférica al radio medio.

\[\fitmath{U=250\approx 250\,\mathrm{m/s}}\]

\item Velocidad axial, constante en el escalón.

\[\fitmath{\mathrm{Cx}=160\approx 160\,\mathrm{m/s}}\]

\item Remolino absoluto a la entrada del rotor.

\[\fitmath{\mathrm{Cth}_{1}=70\approx 70\,\mathrm{m/s}}\]

\item Remolino absoluto a la salida del rotor.

\[\fitmath{\mathrm{Cth}_{2}=180\approx 180\,\mathrm{m/s}}\]

\item Calor específico.

\[\fitmath{\mathrm{cp}=1005\approx 1005\,\mathrm{J/(kg\,K)}}\]

\item Caudal másico.

\[\fitmath{\mathrm{mdot}=20\approx 20\,\mathrm{kg/s}}\]

\item Trabajo específico de Euler con radio constante.

\[\fitmath{w=U\cdot \left(\mathrm{Cth}_{2}-\mathrm{Cth}_{1}\right)}\]

\[\fitmath{w=(250)\cdot \left((180)-(70)\right)\approx 27500\,\mathrm{J/kg}}\]

\item Aumento de temperatura total, adiabático.

\[\fitmath{\mathrm{dT}_{0}=\frac{w}{\mathrm{cp}}}\]

\[\fitmath{\mathrm{dT}_{0}=\frac{(27500)}{(1005)}\approx 27.36318\,\mathrm{K}}\]

\item Temperatura total de salida.

\[\fitmath{T_{02}=T_{01}+\mathrm{dT}_{0}}\]

\[\fitmath{T_{02}=(288.15)+(27.36318408)\approx 315.5132\,\mathrm{K}}\]

\item Rendimiento total-total del escalón.

\[\fitmath{\eta=0.88\approx 0.88\,}\]

\item Razón de presiones total-total.

\[\fitmath{\mathrm{pr}={\left(1+\frac{\eta\cdot \mathrm{dT}_{0}}{T_{01}}\right)}^{\frac{1.4}{0.4}}}\]

\[\fitmath{\mathrm{pr}={\left(1+\frac{(0.88)\cdot (27.36318408)}{(288.15)}\right)}^{\frac{1.4}{0.4}}\approx 1.324323\,}\]

\item Presión total de salida.

\[\fitmath{p_{02}=p_{01}\cdot \mathrm{pr}}\]

\[\fitmath{p_{02}=(101325)\cdot (1.324323484)\approx 134187.1\,\mathrm{Pa}}\]

\item Potencia absorbida por el escalón.

\[\fitmath{P=\mathrm{mdot}\cdot w}\]

\[\fitmath{P=(20)\cdot (27500)\approx 550000\,\mathrm{W}}\]

\item Reacción de un escalón de radio y velocidad axial constantes.

\[\fitmath{R=1-\frac{\mathrm{Cth}_{1}+\mathrm{Cth}_{2}}{2\cdot U}}\]

\[\fitmath{R=1-\frac{(70)+(180)}{2\cdot (250)}\approx 0.5\,}\]

\item Velocidad relativa a la entrada del rotor.

\[\fitmath{W_{1}=\sqrt{{\mathrm{Cx}}^{2}+{\left(U-\mathrm{Cth}_{1}\right)}^{2}}}\]

\[\fitmath{W_{1}=\sqrt{{(160)}^{2}+{\left((250)-(70)\right)}^{2}}\approx 240.8319\,\mathrm{m/s}}\]

\item Velocidad relativa a la salida del rotor.

\[\fitmath{W_{2}=\sqrt{{\mathrm{Cx}}^{2}+{\left(U-\mathrm{Cth}_{2}\right)}^{2}}}\]

\[\fitmath{W_{2}=\sqrt{{(160)}^{2}+{\left((250)-(180)\right)}^{2}}\approx 174.6425\,\mathrm{m/s}}\]

\item Cociente de velocidades relativas (de Haller).

\[\fitmath{\mathrm{DH}=\frac{W_{2}}{W_{1}}}\]

\[\fitmath{\mathrm{DH}=\frac{(174.642492)}{(240.8318916)}\approx 0.7251635\,}\]

\item Ángulo relativo de entrada respecto al eje.

\[\fitmath{b_{1}=\frac{\arctan\left(\frac{U-\mathrm{Cth}_{1}}{\mathrm{Cx}}\right)\cdot 180}{\pi}}\]

\[\fitmath{b_{1}=\frac{\arctan\left(\frac{(250)-(70)}{(160)}\right)\cdot 180}{\pi}\approx 48.36646\,{}^{\circ}}\]

\item Ángulo relativo de salida respecto al eje.

\[\fitmath{b_{2}=\frac{\arctan\left(\frac{U-\mathrm{Cth}_{2}}{\mathrm{Cx}}\right)\cdot 180}{\pi}}\]

\[\fitmath{b_{2}=\frac{\arctan\left(\frac{(250)-(180)}{(160)}\right)\cdot 180}{\pi}\approx 23.62938\,{}^{\circ}}\]

\item Régimen del rodete radial.

\[\fitmath{N=36000\approx 36000\,\mathrm{min^{-1}}}\]

\item Velocidad angular.

\[\fitmath{\mathrm{om}=\frac{2\cdot \pi\cdot N}{60}}\]

\[\fitmath{\mathrm{om}=\frac{2\cdot \pi\cdot (36000)}{60}\approx 3769.911\,\mathrm{rad/s}}\]

\item Velocidad periférica de salida del rodete de 0.30 m de diámetro.

\[\fitmath{U_{2}=\frac{\mathrm{om}\cdot 0.3}{2}}\]

\[\fitmath{U_{2}=\frac{(3769.911184)\cdot 0.3}{2}\approx 565.4867\,\mathrm{m/s}}\]

\item Factor de deslizamiento.

\[\fitmath{\mathrm{sg}=0.85\approx 0.85\,}\]

\item Remolino de salida, sin prerrotación.

\[\fitmath{Cth2r=\mathrm{sg}\cdot U_{2}}\]

\[\fitmath{Cth2r=(0.85)\cdot (565.4866776)\approx 480.6637\,\mathrm{m/s}}\]

\item Trabajo específico de Euler.

\[\fitmath{\mathrm{wr}=U_{2}\cdot Cth2r}\]

\[\fitmath{\mathrm{wr}=(565.4866776)\cdot (480.663676)\approx 271808.9\,\mathrm{J/kg}}\]

\item Temperatura total de entrada del rodete.

\[\fitmath{T01r=300\approx 300\,\mathrm{K}}\]

\item Aumento de temperatura total.

\[\fitmath{dT0r=\frac{\mathrm{wr}}{\mathrm{cp}}}\]

\[\fitmath{dT0r=\frac{(271808.9052)}{(1005)}\approx 270.4566\,\mathrm{K}}\]

\item Rendimiento total-total.

\[\fitmath{\mathrm{etar}=0.8\approx 0.8\,}\]

\item Razón de presiones total-total.

\[\fitmath{\mathrm{prr}={\left(1+\frac{\mathrm{etar}\cdot dT0r}{T01r}\right)}^{\frac{1.4}{0.4}}}\]

\[\fitmath{\mathrm{prr}={\left(1+\frac{(0.8)\cdot (270.4566221)}{(300)}\right)}^{\frac{1.4}{0.4}}\approx 6.689987\,}\]

\item Velocidad del sonido total de entrada.

\[\fitmath{a_{01}=\sqrt{1.4\cdot 287.05287\cdot T01r}}\]

\[\fitmath{a_{01}=\sqrt{1.4\cdot 287.05287\cdot (300)}\approx 347.2207\,\mathrm{m/s}}\]

\item Mach periférico de salida.

\[\fitmath{\mathrm{Mu}=\frac{U_{2}}{a_{01}}}\]

\[\fitmath{\mathrm{Mu}=\frac{(565.4866776)}{(347.2206869)}\approx 1.628609\,}\]

\item Potencia con 1.2 kg/s.

\[\fitmath{\mathrm{Pr}=1.2\cdot \mathrm{wr}}\]

\[\fitmath{\mathrm{Pr}=1.2\cdot (271808.9052)\approx 326170.7\,\mathrm{W}}\]

\end{enumerate}


\section{Simulacro D}
\textbf{Cuestiones.} El coeficiente de carga $\psi=w/U^2$ mide el trabajo extraído por unidad de energía cinética periférica y el de flujo $\phi=C_x/U$ el caudal axial adimensional; juntos fijan los ángulos de los triángulos y, con ellos, la difusión y las pérdidas. La reacción es el cociente de la caída de entalpía estática del rotor entre la del escalón; con radio constante y rotor adiabático la caída estática del rotor es $(W_3^2-W_2^2)/2$, y en un escalón repetitivo ($C_3=C_1$) la caída del escalón coincide con el trabajo. En una barra uniforme la tensión en la raíz es el peso centrífugo del material por encima de ella dividido por el área, y el área se cancela.

\textbf{Escalón.} $w=UC_{\theta2}=122880\,\mathrm{J/kg}$, $\psi=1.2$, $\phi=0.625$ y $\Delta T_0=106.852\,\mathrm K$. La razón de expansión es $1.42925$ y el ángulo absoluto de entrada al rotor es $62.488^\circ$. Las velocidades relativas al cuadrado son $W_2^2=44096\,\mathrm{m^2/s^2}$ y $W_3^2=142400\,\mathrm{m^2/s^2}$, y la reacción es $(W_3^2-W_2^2)/(2w)=0.4$.

\textbf{Pala.} $\sigma=\rho\omega^2(r_t^2-r_h^2)/2=2.91351\times10^{8}\,\mathrm{Pa}$ con $U_t=424.115\,\mathrm{m/s}$, de modo que el coeficiente de seguridad frente a $400\,\mathrm{MPa}$ es $1.37292$. La red térmica tiene $R=0.00101556\,\mathrm{m^2K/W}$, un flujo $q=787746\,\mathrm{W/m^2}$, una pared caliente a $1062.36\,\mathrm K$ y una diferencia de temperatura de $47.2648\,\mathrm K$ a través de la pared. La tensión térmica con restricción total es $E\alpha\Delta T=1.25724\times10^{8}\,\mathrm{Pa}$. El modelo deja fuera la concentración de tensiones en la raíz, la fluencia a esa temperatura, la fatiga, la película de refrigeración y el gradiente térmico bidimensional en el borde de ataque.

\subsection*{Procedimiento y operaciones del simulacro D}
El escalón repetitivo con descarga axial se resuelve con Euler; la reacción sale de la caída estática del rotor entre el trabajo, con radio constante. La pala se trata como barra uniforme sometida a fuerza centrífuga y la pared como una red térmica de tres resistencias en serie.

\begin{enumerate}

\item Temperatura total de entrada a la turbina.

\[\fitmath{T_{01}=1400\approx 1400\,\mathrm{K}}\]

\item Calor específico del gas caliente.

\[\fitmath{\mathrm{cp}=1150\approx 1150\,\mathrm{J/(kg\,K)}}\]

\item Razón de calores del gas caliente.

\[\fitmath{\gamma=1.33\approx 1.33\,}\]

\item Velocidad periférica al radio medio.

\[\fitmath{U=320\approx 320\,\mathrm{m/s}}\]

\item Velocidad axial constante.

\[\fitmath{\mathrm{Cx}=200\approx 200\,\mathrm{m/s}}\]

\item Remolino absoluto a la salida del estator.

\[\fitmath{\mathrm{Cth}_{2}=384\approx 384\,\mathrm{m/s}}\]

\item Trabajo específico con descarga axial.

\[\fitmath{w=U\cdot \mathrm{Cth}_{2}}\]

\[\fitmath{w=(320)\cdot (384)\approx 122880\,\mathrm{J/kg}}\]

\item Coeficiente de carga.

\[\fitmath{\psi=\frac{w}{{U}^{2}}}\]

\[\fitmath{\psi=\frac{(122880)}{{(320)}^{2}}\approx 1.2\,}\]

\item Coeficiente de flujo.

\[\fitmath{\varphi=\frac{\mathrm{Cx}}{U}}\]

\[\fitmath{\varphi=\frac{(200)}{(320)}\approx 0.625\,}\]

\item Caída de temperatura total.

\[\fitmath{\mathrm{dT}_{0}=\frac{w}{\mathrm{cp}}}\]

\[\fitmath{\mathrm{dT}_{0}=\frac{(122880)}{(1150)}\approx 106.8522\,\mathrm{K}}\]

\item Rendimiento total-total.

\[\fitmath{\eta=0.9\approx 0.9\,}\]

\item Razón de expansión total-total.

\[\fitmath{\mathrm{pr}={\left(1-\frac{\mathrm{dT}_{0}}{\eta\cdot T_{01}}\right)}^{\frac{-\gamma}{\gamma-1}}}\]

\[\fitmath{\mathrm{pr}={\left(1-\frac{(106.8521739)}{(0.9)\cdot (1400)}\right)}^{\frac{-(1.33)}{(1.33)-1}}\approx 1.429251\,}\]

\item Ángulo absoluto a la entrada del rotor.

\[\fitmath{a_{2}=\frac{\arctan\left(\frac{\mathrm{Cth}_{2}}{\mathrm{Cx}}\right)\cdot 180}{\pi}}\]

\[\fitmath{a_{2}=\frac{\arctan\left(\frac{(384)}{(200)}\right)\cdot 180}{\pi}\approx 62.488\,{}^{\circ}}\]

\item Cuadrado de la velocidad relativa a la entrada del rotor.

\[\fitmath{W2sq={\mathrm{Cx}}^{2}+{\left(\mathrm{Cth}_{2}-U\right)}^{2}}\]

\[\fitmath{W2sq={(200)}^{2}+{\left((384)-(320)\right)}^{2}\approx 44096\,\mathrm{m^{2}/s^{2}}}\]

\item Cuadrado de la velocidad relativa a la salida, con descarga axial.

\[\fitmath{W3sq={\mathrm{Cx}}^{2}+{U}^{2}}\]

\[\fitmath{W3sq={(200)}^{2}+{(320)}^{2}\approx 142400\,\mathrm{m^{2}/s^{2}}}\]

\item Reacción como cociente de caída estática del rotor y trabajo del escalón repetitivo.

\[\fitmath{R=\frac{W3sq-W2sq}{2\cdot w}}\]

\[\fitmath{R=\frac{(142400)-(44096)}{2\cdot (122880)}\approx 0.4\,}\]

\item Densidad del material de la pala.

\[\fitmath{\rho=8200\approx 8200\,\mathrm{kg/m^{3}}}\]

\item Régimen de giro.

\[\fitmath{\mathrm{Nb}=9000\approx 9000\,\mathrm{min^{-1}}}\]

\item Velocidad angular.

\[\fitmath{\mathrm{om}=\frac{2\cdot \pi\cdot \mathrm{Nb}}{60}}\]

\[\fitmath{\mathrm{om}=\frac{2\cdot \pi\cdot (9000)}{60}\approx 942.4778\,\mathrm{rad/s}}\]

\item Tensión centrífuga de una barra uniforme entre cubo y punta.

\[\fitmath{\mathrm{sig}=\frac{\rho\cdot {\mathrm{om}}^{2}\cdot \left({0.45}^{2}-{0.35}^{2}\right)}{2}}\]

\[\fitmath{\mathrm{sig}=\frac{(8200)\cdot {(942.4777961)}^{2}\cdot \left({0.45}^{2}-{0.35}^{2}\right)}{2}\approx 2.913507\times10^{8}\,\mathrm{Pa}}\]

\item Velocidad periférica en la punta.

\[\fitmath{\mathrm{Ut}=\mathrm{om}\cdot 0.45}\]

\[\fitmath{\mathrm{Ut}=(942.4777961)\cdot 0.45\approx 424.115\,\mathrm{m/s}}\]

\item Tensión admisible docente a la temperatura de la pala.

\[\fitmath{\mathrm{sa}=400000000\approx 4\times10^{8}\,\mathrm{Pa}}\]

\item Coeficiente de seguridad.

\[\fitmath{\mathrm{SF}=\frac{\mathrm{sa}}{\mathrm{sig}}}\]

\[\fitmath{\mathrm{SF}=\frac{(400000000)}{(291350721.9)}\approx 1.372916\,}\]

\item Temperatura de recuperación exterior.

\[\fitmath{\mathrm{Tg}=1500\approx 1500\,\mathrm{K}}\]

\item Temperatura del refrigerante.

\[\fitmath{\mathrm{Tc}=700\approx 700\,\mathrm{K}}\]

\item Resistencia térmica por área de la red de pared.

\[\fitmath{\mathrm{Rt}=\frac{1}{1800}+\frac{0.0015}{25}+\frac{1}{2500}\approx 0.001015556\,\mathrm{m^{2}\,K/W}}\]

\item Flujo de calor estacionario.

\[\fitmath{q=\frac{\mathrm{Tg}-\mathrm{Tc}}{\mathrm{Rt}}}\]

\[\fitmath{q=\frac{(1500)-(700)}{(0.001015555556)}\approx 787746.2\,\mathrm{W/m^{2}}}\]

\item Temperatura de la pared caliente.

\[\fitmath{\mathrm{Twh}=\mathrm{Tg}-\frac{q}{1800}}\]

\[\fitmath{\mathrm{Twh}=(1500)-\frac{(787746.1707)}{1800}\approx 1062.363\,\mathrm{K}}\]

\item Diferencia de temperatura a través de la pared.

\[\fitmath{\mathrm{dTw}=\frac{q\cdot 0.0015}{25}}\]

\[\fitmath{\mathrm{dTw}=\frac{(787746.1707)\cdot 0.0015}{25}\approx 47.26477\,\mathrm{K}}\]

\item Tensión térmica con restricción total uniaxial.

\[\fitmath{\mathrm{sth}=1.9\times10^{11}\cdot 1.4\times10^{-5}\cdot \mathrm{dTw}}\]

\[\fitmath{\mathrm{sth}=1.9\times10^{11}\cdot 1.4\times10^{-5}\cdot (47.26477024)\approx 1.257243\times10^{8}\,\mathrm{Pa}}\]

\end{enumerate}


\section{Simulacro E}
\textbf{Cuestiones.} El régimen corregido $N/\sqrt\theta$ y el caudal corregido $\dot m\sqrt\theta/\delta$ agrupan los efectos de la temperatura y la presión de entrada si el gas es el mismo, el número de Reynolds es comparable y la cinemática es semejante; entonces la razón de presiones y el rendimiento son iguales. Un margen de bombeo es una distancia a la línea de bombeo del mapa, y esa línea solo se conoce con datos experimentales del compresor completo. La turbina entrega la potencia que absorbe el compresor y los auxiliares, más las pérdidas mecánicas, porque comparten eje: si hubiera un exceso, el régimen aumentaría.

\textbf{Corrección.} $\theta=0.867604$ y $\delta=0.690846$. El régimen corregido es $10950.6\,\mathrm{min^{-1}}$ y el caudal corregido $18.8759\,\mathrm{kg/s}$. Para $T=320\,\mathrm K$ y $p=101325\,\mathrm{Pa}$, $\theta=1.11053$ y $\delta=1$: el régimen físico es $11540\,\mathrm{min^{-1}}$ y el caudal físico $17.9119\,\mathrm{kg/s}$. Se supone el mismo gas, número de Reynolds comparable y la misma cinemática; no se corrige por holguras ni por cambios de $\gamma$.

\textbf{Generador de gas.} La potencia del compresor es $5.3064\times10^{6}\,\mathrm W$. El caudal de turbina es $21.34\,\mathrm{kg/s}$ y la caída de temperatura total necesaria, $\Delta T_{0t}=P_c/(\eta_m\dot m_tc_{p})=218.41\,\mathrm K$. Con $\eta_t=0.88$ y $T_{04}=1450\,\mathrm K$ la razón de expansión total-total es $2.1311$. El modelo omite el aire de refrigeración reinyectado, la variación de $c_p$ con la temperatura y la potencia de los auxiliares.

\subsection*{Procedimiento y operaciones del simulacro E}
La corrección usa los cocientes $\theta=T/T_{ref}$ y $\delta=p/p_{ref}$: $N/\sqrt\theta$ y $\dot m\sqrt\theta/\delta$. El acoplamiento iguala la potencia del compresor con la de la turbina multiplicada por el rendimiento mecánico, y la razón de expansión sale de la caída de temperatura con el rendimiento total-total.

\begin{enumerate}

\item Temperatura total de entrada en el ensayo.

\[\fitmath{\mathrm{Tin}=250\approx 250\,\mathrm{K}}\]

\item Presión total de entrada en el ensayo.

\[\fitmath{\mathrm{pin}=70000\approx 70000\,\mathrm{Pa}}\]

\item Caudal másico medido.

\[\fitmath{m=14\approx 14\,\mathrm{kg/s}}\]

\item Régimen medido.

\[\fitmath{N=10200\approx 10200\,\mathrm{min^{-1}}}\]

\item Cociente de temperatura respecto a la referencia.

\[\fitmath{\mathrm{th}=\frac{\mathrm{Tin}}{288.15}}\]

\[\fitmath{\mathrm{th}=\frac{(250)}{288.15}\approx 0.8676037\,}\]

\item Cociente de presión respecto a la referencia.

\[\fitmath{\mathrm{de}=\frac{\mathrm{pin}}{101325}}\]

\[\fitmath{\mathrm{de}=\frac{(70000)}{101325}\approx 0.6908463\,}\]

\item Régimen corregido.

\[\fitmath{\mathrm{Nc}=\frac{N}{\sqrt{\mathrm{th}}}}\]

\[\fitmath{\mathrm{Nc}=\frac{(10200)}{\sqrt{(0.8676036786)}}\approx 10950.64\,\mathrm{min^{-1}}}\]

\item Caudal corregido.

\[\fitmath{\mathrm{mc}=\frac{m\cdot \sqrt{\mathrm{th}}}{\mathrm{de}}}\]

\[\fitmath{\mathrm{mc}=\frac{(14)\cdot \sqrt{(0.8676036786)}}{(0.6908462867)}\approx 18.87588\,\mathrm{kg/s}}\]

\item Cociente de temperatura del punto equivalente.

\[\fitmath{\mathrm{th}_{2}=\frac{320}{288.15}\approx 1.110533\,}\]

\item Régimen físico con el mismo régimen corregido.

\[\fitmath{N_{2}=\mathrm{Nc}\cdot \sqrt{\mathrm{th}_{2}}}\]

\[\fitmath{N_{2}=(10950.63943)\cdot \sqrt{(1.110532709)}\approx 11539.98\,\mathrm{min^{-1}}}\]

\item Caudal físico a 101325 Pa con el mismo caudal corregido.

\[\fitmath{m_{2}=\frac{\mathrm{mc}}{\sqrt{\mathrm{th}_{2}}}}\]

\[\fitmath{m_{2}=\frac{(18.87588403)}{\sqrt{(1.110532709)}}\approx 17.9119\,\mathrm{kg/s}}\]

\item Caudal del compresor.

\[\fitmath{\mathrm{mcomp}=22\approx 22\,\mathrm{kg/s}}\]

\item Aumento de temperatura total del compresor.

\[\fitmath{\mathrm{dTc}=240\approx 240\,\mathrm{K}}\]

\item Potencia absorbida por el compresor.

\[\fitmath{\mathrm{Pc}=\mathrm{mcomp}\cdot 1005\cdot \mathrm{dTc}}\]

\[\fitmath{\mathrm{Pc}=(22)\cdot 1005\cdot (240)\approx 5306400\,\mathrm{W}}\]

\item Caudal de la turbina con combustible y sangrado.

\[\fitmath{\mathrm{mt}=\left(1+0.02-0.05\right)\cdot \mathrm{mcomp}}\]

\[\fitmath{\mathrm{mt}=\left(1+0.02-0.05\right)\cdot (22)\approx 21.34\,\mathrm{kg/s}}\]

\item Caída de temperatura total requerida en la turbina.

\[\fitmath{\mathrm{dTt}=\frac{\mathrm{Pc}}{0.99\cdot \mathrm{mt}\cdot 1150}}\]

\[\fitmath{\mathrm{dTt}=\frac{(5306400)}{0.99\cdot (21.34)\cdot 1150}\approx 218.41\,\mathrm{K}}\]

\item Temperatura total de entrada a la turbina.

\[\fitmath{T_{04}=1450\approx 1450\,\mathrm{K}}\]

\item Razón de expansión total-total de la turbina.

\[\fitmath{\mathrm{prt}={\left(1-\frac{\mathrm{dTt}}{0.88\cdot T_{04}}\right)}^{\frac{-1.33}{0.33}}}\]

\[\fitmath{\mathrm{prt}={\left(1-\frac{(218.4100077)}{0.88\cdot (1450)}\right)}^{\frac{-1.33}{0.33}}\approx 2.131101\,}\]

\end{enumerate}

